\section{Closed boundary cases and general reductions}

\subsection{Commuting triples}
\begin{proposition}[Commuting case \statusR]
If $X,Y,Z$ commute, then \eqref{eq:MY} holds for both $j=1,2$.
\end{proposition}
\begin{proof}
The three matrices are simultaneously diagonalizable.  On each common eigenline the desired estimate becomes the scalar weighted Young inequality
\[
x^\alpha y^\beta z^\gamma\le \alpha x+\beta y+\gamma z.
\]
Ordering the two resulting scalar coordinates gives the required eigenvalue inequalities.
\end{proof}

\subsection{One scalar factor}
\begin{proposition}[Scalar-factor case \statusR]
If one of $X,Y,Z$ is a scalar matrix, then \eqref{eq:MY} holds in every dimension.
\end{proposition}
\begin{proof}
It suffices to take $Z=zI$.  Set
\[
\delta=\alpha+\beta,\qquad a=\frac{\alpha}{\delta},\qquad b=\frac{\beta}{\delta}.
\]
Then $a+b=1$ and
\[
X^\alpha Y^\beta=(X^\delta)^a(Y^\delta)^b.
\]
By Ando's two-factor Matrix Young inequality,
\[
s_j(X^\alpha Y^\beta)
\le \lambda_j\bigl(aX^\delta+bY^\delta\bigr).
\]
Because $0<\delta<1$, operator concavity of $t\mapsto t^\delta$ gives
\[
aX^\delta+bY^\delta\preceq (aX+bY)^\delta.
\]
Functional calculus and the scalar weighted Young inequality imply
\[
z^\gamma(aX+bY)^\delta
\preceq \delta(aX+bY)+\gamma zI
=\alpha X+\beta Y+\gamma zI.
\]
Combining the inequalities proves the claim.
\end{proof}

\subsection{Adjacent identical factors}
\begin{proposition}[Adjacent-equal case \statusR]
If $X=Y$ or $Y=Z$, then \eqref{eq:MY} reduces to the two-factor Ando inequality.
\end{proposition}
\begin{proof}
For example, if $X=Y$ then
\[
X^\alpha Y^\beta Z^\gamma=X^{\alpha+\beta}Z^\gamma,
\]
and $(\alpha+\beta)+\gamma=1$.
\end{proof}

\subsection{A determinant slack available in dimension two}
For the remainder of the manuscript let $n=2$ and assume positive definiteness when inverses or logarithms are used; semidefinite cases can be obtained by regularization.

Put
\[
P=X^\alpha Y^\beta Z^\gamma,
\qquad
S=\alpha X+\beta Y+\gamma Z.
\]
Define
\[
\rho=\sqrt{\det P},\qquad \Delta=\sqrt{\det S}.
\]
Since $\det^{1/2}$ is concave on $\mathbb H_+^2$,
\[
\Delta\ge \alpha\sqrt{\det X}+\beta\sqrt{\det Y}+\gamma\sqrt{\det Z}.
\]
Applying scalar weighted AM--GM to the three positive numbers $\sqrt{\det X}$, $\sqrt{\det Y}$ and $\sqrt{\det Z}$ yields
\begin{equation}
\boxed{\Delta\ge \rho.}
\label{eq:det-slack}
\end{equation}
This nonnegative determinant slack is one of the principal special features of the $2\times2$ problem.
